\subsection{Contribution of Copper Mining to Global Emissions}

Copper plays a central role in the global energy transition due to its extensive use in renewable energy systems, electric vehicles, electricity grids, and energy storage technologies. According to the International Energy Agency (IEA), electric vehicles require approximately 2.5 times more copper than conventional internal combustion vehicles \cite{IEA2021}. As a result, several studies project a significant increase in copper demand under global decarbonization scenarios \cite{Watari2020}. The IEA estimates that copper demand from clean energy technologies alone could more than double by 2040 under a net-zero emissions pathway \cite{IEA2021}.

At the same time, copper extraction and processing remain highly energy-intensive activities requiring large quantities of fuel and electricity \cite{Norgate2010}. Chile currently accounts for more than 25\% of global copper production, making the environmental performance of its mining sector particularly important at the global scale \cite{CochilcoWebsite2025}. Consequently, a growing body of literature focuses on the environmental sustainability of copper supply chains and the challenge of reducing emissions associated with mining and mineral processing activities \cite{Elshkaki2018}.


\subsection{Historical Evolution of Mining Energy Consumption}
\label{sec:evol of mining energy}

Several studies have examined the historical evolution of energy consumption in metallic mining and its relationship with industrial development \cite{Norgate2010,Northey2014}. In copper mining, energy demand has increased steadily over recent decades alongside the expansion and increasing complexity of mining operations.

Mining energy consumption is generally divided into direct and indirect components. Direct energy consumption mainly corresponds to fuel combustion in trucks, excavation equipment, and transportation systems, while indirect energy consumption refers to electricity use in crushing, grinding, concentration, pumping, and refining processes. Existing studies emphasize that electricity consumption has become increasingly important in modern copper mining because mineral processing operations are highly electricity-intensive \cite{Norgate2010}. For example, concentrator plants in Chile consume more than 10,000 MJ of electricity per tonne of copper concentrate produced \cite{Cochilco2024}.

The literature also highlights that technological improvements and process optimization have partially reduced energy requirements in certain mining activities. However, total energy consumption in the sector has generally continued to increase due to the scale and growing complexity of production systems \cite{Calvo2016}.


\subsection{Ore Grade Decline and Energy Demand}

Declining ore grades are widely recognized as one of the main structural challenges facing the mining industry \cite{Calvo2016,Northey2014}. Ore grade refers to the concentration of valuable metal contained in the extracted rock. As ore grades decline, larger quantities of material must be extracted and processed in order to produce the same quantity of refined copper.

Several studies identify ore grade decline as a major driver of increasing energy consumption in copper mining \cite{Calvo2016}. Lower ore concentrations require additional crushing, grinding, concentration, and waste management operations, all of which increase energy requirements throughout the production chain. Northey et al. show that average copper ore grades in many mining operations have fallen below 1\% over recent decades \cite{Northey2014}. In Chile, average ore grades in concentrator plants declined from approximately 0.87\% in 2015 to around 0.66\% in 2024 \cite{Cochilco2024}.

Because ore grade is determined by geological conditions and resource depletion, this trend represents a long-term structural constraint that may progressively increase the energy requirements associated with copper production.


\subsection{Electrification and Renewable Integration}
\label{sec:elec & renew}
In response to growing environmental concerns, many studies explore potential decarbonization pathways for the mining industry. Electrification of mining equipment and integration of renewable electricity are among the most frequently proposed strategies for reducing greenhouse gas emissions in copper production \cite{Memary2022}. Existing research suggests that replacing diesel-powered systems with electric technologies could reduce direct emissions from mining operations, while decarbonizing electricity generation could lower indirect emissions associated with electricity-intensive mineral processing.

Several studies identify Chile as particularly well positioned for this transition because of its exceptional renewable energy potential, especially solar energy in the Atacama Desert and wind energy in coastal regions \cite{IEAChile2026}. Renewable sources represented nearly 70\% of Chile’s electricity generation in 2024, while wind and solar energy alone exceeded 40\% of monthly electricity generation for the first time in December 2024 \cite{EnergiaEstrategica2025,Ember2025}.

The literature nevertheless emphasizes that reducing emissions from copper mining will likely require a combination of renewable electricity integration, technological innovation, process efficiency improvements, and broader industrial decarbonization strategies \cite{Memary2022}.


\subsection{Methodologies for Emissions Decomposition}

Several methodological approaches have been developed to analyse greenhouse gas emissions in industrial sectors, including decomposition methods such as the Kaya identity, IPAT frameworks, and Logarithmic Mean Divisia Index (LMDI) techniques. These approaches are widely used to separate the influence of activity growth, energy intensity, and carbon intensity on total emissions \cite{Tilton1998,Norgate2010}.

In the context of copper mining, decomposition methods are particularly relevant because emissions are influenced simultaneously by production growth, ore quality, energy consumption, and the carbon intensity of fuel and electricity use. These interacting drivers motivate the use of an IPAT-based framework capable of separating geological, technological, and energy-related contributions to CO$_2$ emissions.

